\section{What existing sequence indexes store}
\label{sec:index-minimap2-comparison}

\subsection{Minimap2: buckets, singletons, and position lists}

Minimap2 selects minimizers from the reference and groups their occurrences
by hash\autocite{minimap2_paper,minimap2_software}. Its index construction
first uses the lowest $b$ hash bits to choose one of $2^b$ buckets. Within
each bucket, sorting places equal hashes together. A bucket-local hash table
then records where their positions can be found. The construction issue and
the accompanying format explanation give a useful view of this
layout\autocite{minimap2_issue_820,minimap2_mmi_gist}.

There are two cases:
\begin{itemize}
  \item A \textbf{singleton} occurs once. Its position is stored directly
    in the hash-table value.
  \item A \textbf{repeated key} has several occurrences. Its value stores
    an offset and a count into the bucket's separate position array.
\end{itemize}
Each position record includes the reference sequence identifier, position,
and strand. A marker bit in the stored key distinguishes the two cases.
Figure~\ref{fig:index-mmi-construction} uses small integer hashes to show the
partition and lookup; it omits the marker's bit encoding. Its positions are
zero-based last-symbol coordinates.

\begin{figure}[htbp]
\centering
\begin{tikzpicture}[x=1cm,y=1cm,font=\footnotesize]
  \node[anchor=west,font=\bfseries] at (0,5.8) {Minimizer occurrences before grouping};
  \foreach \x/\h/\p in {1/13/17,3/21/24,5/10/29,7/21/30}
    \node[index blue,minimum width=1.6cm] at (\x,4.95) {$h=\h$\\$p=\p$};
  \node[anchor=west,align=left] at (9,4.95) {$b=2$: bucket $h\bmod4$\\remaining key: $\lfloor h/4\rfloor$};
  \foreach \x/\b in {1/0,4/1,7/2,10/3}
    \node[index cell,fill=black!3,minimum width=2.2cm] (b\b) at (\x,3.3) {bucket \b};
  \draw[index arrow,indexBlue] (1,4.42)--(4,3.6);
  \draw[index arrow,indexBlue] (3,4.42)--(4,3.6);
  \draw[index arrow,indexBlue] (7,4.42)--(4,3.6);
  \draw[index arrow,black!45] (5,4.42)--(7,3.6);
  \node[anchor=west,font=\bfseries] at (0,2.35) {Bucket 1 after sorting};
  \node[index blue,minimum width=3.6cm] (single) at (2.35,1.5)
    {remaining bits 3: $p=17$};
  \node[index orange,minimum width=3.6cm] (repeated) at (2.35,.5)
    {remaining bits 5: offset 0, count 2};
  \node[anchor=west] at (6.0,1.7) {separate position array};
  \node[index orange,minimum width=1.15cm] at (6.65,.5) {$p=24$};
  \node[index orange,minimum width=1.15cm] at (7.85,.5) {$p=30$};
  \node[index cell,dashed,minimum width=1.15cm] at (9.05,.5) {$\cdots$};
  \node at (6.65,1.08) {0}; \node at (7.85,1.08) {1};
  \draw[index arrow,indexOrange] (repeated.east)--(6.05,.5);
  \node[anchor=west,text=black!65] at (0,-.45)
    {A singleton needs no separate list. Repeated keys share a contiguous array.};
\end{tikzpicture}
\caption{A simplified minimap2 bucket. Hashes 13 and 21 have the same
lowest two bits, but different remaining keys. The two occurrences of hash
21 become adjacent in the bucket's position array. Concept based on the
minimap2 implementation and format discussion.}
\label{fig:index-mmi-construction}
\end{figure}

The \texttt{.mmi} writer saves parameters, reference names and lengths,
bucket position arrays, hash entries, and optionally packed reference symbols.
The loader allocates the bucket structures and reconstructs one complete index
part. A large reference can be split into several parts, so ``complete'' here
does not require every part to be resident at once.

\noindent
\coderef{minimap2index/index.c}{Minimizer bucket construction}{worker_post}{226}

\noindent
\coderef{minimap2index/index.c}{Writing and loading an index part}{mm_idx_dump, mm_idx_load}{475}

This is a useful arrangement for a mapper whose caller chooses a reference
file. BioSeqDB keeps the same useful key-to-occurrences idea, but adds
selective database reads and source-version checks. These are concrete
differences in responsibility, not a measured speed comparison.

\subsection{MMseqs2: words point back to target sequences}

The CPU MMseqs2 prefilter uses an inverted index: a word points to the target
sequences and positions containing it. Its 2017 description uses
\emph{spaced k-mers}, words selected through a fixed pattern that skips some
positions. Query processing also generates similar words according to protein
substitution scores, instead of requiring every matched word to be
identical\autocite{Steinegger079681,mmseqs2_2017}.

The lookup has three steps:
\begin{enumerate}
  \item Take a spaced word from the query and generate similar words.
  \item Read their posting lists, whose entries are pairs
    $(\text{target identifier},\text{target position})$.
  \item Find successive matches to one target on the same \emph{diagonal}, defined as
    target position minus query position. These matches suggest an ungapped
    candidate region for further scoring.
\end{enumerate}

The current ordinary-sequence builder counts entries for each encoded word,
forms cumulative offsets, and fills one contiguous entry array. Looking up
code $c$ then reads entries in $[offset[c],offset[c+1])$. Offsets count
entries, not bytes. The offset table reserves a slot for every possible
encoded word, including absent words. This makes a list easy to find, but
the table's size grows with the alphabet and word length.
The implementation keeps the first position per
distinct word code in each target sequence; it does not promise every
repeated occurrence within that sequence. BioSeqDB's minimizer indexes
instead retain distinct indexed minimizer occurrences, including occurrences
at different positions in one target. These positions can support different
alignments from the same stored row.

\noindent
\coderef{mmseqsindex/src/prefiltering/IndexBuilder.cpp}{Count and fill the MMseqs2 index}{IndexBuilder::fillDatabase}{56}

\noindent
\coderef{mmseqsindex/src/prefiltering/IndexTable.h}{Read a word's entry range}{getDBSeqList}{178}

\noindent
\coderef{mmseqsindex/src/prefiltering/IndexTable.h}{Keep distinct word codes within one target}{addSequence}{350}

\begin{figure}[htbp]
\centering
\begin{tikzpicture}[x=1cm,y=1cm,font=\footnotesize]
  \node[anchor=west,font=\bfseries] at (0,6.55) {(a) Read target words};
  \foreach \x/\id/\s in {2/1/ACD,6/2/CAD,10/3/ACE}
    \node[index blue,minimum width=2.8cm] at (\x,5.75) {target \id: \texttt{\s}};
  \node[anchor=west,font=\bfseries] at (0,4.8) {(b) Count entries, then sum counts to find list starts};
  \node[anchor=east] at (1.25,3.95) {word};
  \node[anchor=east] at (1.25,3.2) {count};
  \node[anchor=east] at (1.25,2.45) {offset};
  \foreach \x/\word/\n/\o in {2.25/AC/2/0,4.1/AD/1/2,5.95/CA/1/3,7.8/CD/1/4,9.65/CE/1/5} {
    \node[index blue,minimum width=1.55cm] at (\x,3.95) {\texttt{\word}};
    \node[index cell,minimum width=1.55cm] at (\x,3.2) {$\n$};
    \node[index green,minimum width=1.55cm] at (\x,2.45) {$\o$};
  }
  \node[index green,minimum width=1.25cm] at (11.25,2.45) {6};
  \node[above] at (11.25,2.78) {end};
  \node[anchor=west,font=\bfseries] at (0,1.5) {(c) Fill the reserved list ranges};
  \foreach \x/\id/\p/\i in {2.1/1/0/0,3.7/3/0/1,5.3/2/1/2,6.9/2/0/3,8.5/1/1/4,10.1/3/1/5} {
    \node[index cell,minimum width=1.48cm] at (\x,.5) {$(\id,\p)$};
    \node[above] at (\x,.81) {$\i$};
  }
  \draw[indexBlue,very thick] (1.3,.12)--(4.5,.12);
  \node[below,text=indexBlue] at (2.9,.12) {AC: entries $[0,2)$};
  \node[anchor=west,align=left] at (6,-.15)
    {Entry = (target ID, word start).\\Positions start at zero.};
\end{tikzpicture}
\caption{A small example of MMseqs2's count-and-fill construction, using
contiguous two-letter words for clarity. The real CPU search uses its
configured spaced-word pattern and larger words. Counts reserve adjacent
list ranges, and a second pass fills them. AC occurs in targets 1 and 3.
Empty code slots are omitted from the drawing.}
\label{fig:index-mmseqs-inverted}
\end{figure}

MMseqs2's later tools demonstrate why the distinction matters. Linclust
selects a limited number of low-hash words per sequence and sorts a table
containing the word, sequence identifier, and position. Equal words bring
candidate sequences together for clustering\autocite{mmseqs2_linclust_2018}.
This whole-sequence sampling differs from choosing a minimizer in each
sliding window. The interactive server precomputes its CPU index and other
lookup structures, then uses memory mapping so repeated searches can benefit
from the operating system's page cache\autocite{mmseqs2_desktop_2019}.

The same software family also serves taxonomy searches, where sequence hits
must remain connected to taxonomic identifiers and reference
metadata\autocite{mmseqs2_taxonomy_2021}. Its GPU search follows another
route: it scans sequence residues with gapless dynamic programming, then
performs selected gapped alignments. It can keep database batches on the GPU
or stream them as needed\autocite{mmseqs2_gpu_2025}. Thus a persistent
word index, a sampled clustering table, and resident sequence batches solve
different problems. BioSeqDB's persistent minimizer indexes serve its
seed-and-chain alignment path.
